\documentclass[11pt,letterpaper]{article}
\usepackage[T1]{fontenc}
\usepackage[utf8]{inputenc}
\usepackage{lmodern}
\usepackage[margin=0.95in]{geometry}
\usepackage{amsmath,amssymb,amsthm}
\usepackage{enumitem}
\usepackage{microtype}
\usepackage{xcolor}
\usepackage{booktabs,array}
\usepackage{fancyhdr}
\usepackage{hyperref}
\definecolor{mathjockgreen}{HTML}{355439}
\hypersetup{colorlinks=true,linkcolor=mathjockgreen,urlcolor=mathjockgreen,
  pdftitle={Ross Elementary Analysis: Chapter 1 worked solutions},
  pdfauthor={Mathjock},pdfsubject={Detailed Chapter 1 solutions, second edition}}
\setlength{\parindent}{0pt}
\setlength{\parskip}{0.55em}
\setlength{\headheight}{14pt}
\setlist[enumerate]{leftmargin=*,itemsep=0.4em,topsep=0.4em}
\setlist[itemize]{leftmargin=*,itemsep=0.25em}
\allowdisplaybreaks
\emergencystretch=2em
\pagestyle{fancy}
\fancyhf{}
\fancyhead[L]{\small Mathjock \textbar\ Ross, Chapter 1}
\fancyhead[R]{\small Step-by-step solutions}
\fancyfoot[C]{\thepage}
\newcommand{\N}{\mathbb{N}}
\newcommand{\Z}{\mathbb{Z}}
\newcommand{\Q}{\mathbb{Q}}
\newcommand{\R}{\mathbb{R}}
\newcommand{\abs}[1]{\left|#1\right|}
\newcommand{\exercise}[2]{\subsection{Exercise #1: #2}\label{ex:#1}}
\newcommand{\goal}{\par\textbf{What we must show.} }
\newcommand{\idea}{\par\textbf{Plan.} }
\newcommand{\step}[1]{\par\medskip\textbf{#1}\par}
\newcommand{\answer}[1]{\par\textbf{Conclusion.} #1\par}
\newcommand{\bookinfo}{%
\textbf{Text and numbering.} Kenneth A. Ross, \emph{Elementary Analysis:
The Theory of Calculus}, second edition (Springer, 2013), Chapter~1.
Exercise numbers follow the copy already linked on Mathjock. These are
independently written worked solutions, not the publisher's solutions manual.
Each exercise is identified by a short mathematical description.\par
\textbf{Reading philosophy.} Every argument states its goal, records the
assumptions being used, and explains the intermediate algebra or logical
inference. Definitions, sign conditions, and both directions of an
``if and only if'' are made explicit. A cited rule is named at the point
where it is used.\par
\textbf{Conventions.} $\N=\{1,2,3,\ldots\}$,
$\Z=\{\ldots,-2,-1,0,1,2,\ldots\}$, and
$\Q=\{p/q:p,q\in\Z,\ q\ne0\}$. Unless a different domain is stated,
letters in Section~3 denote real numbers. The empty sum is $0$, and $0!=1$.
}
